\section{Validation statistique}
\label{sec:validation}

\subsection{Validation croisée}

Une validation croisée à 5 plis a été réalisée sur l'ensemble d'entraînement pour
le meilleur pipeline (Full(opt) + SVM) :

\begin{table}[H]
  \centering
  \caption{Résultats de la validation croisée à 5 plis pour Full(opt) + SVM.}
  \label{tab:cv}
  \begin{tabular}{@{}lcc@{}}
    \toprule
    \textbf{Métrique} & \textbf{Moyenne} & \textbf{Écart-type} \\
    \midrule
    F1-macro & \BestFOneCVMean & \BestFOneCVStd \\
    Précision & \BestAccCVMean  & \BestAccCVStd \\
    \bottomrule
  \end{tabular}
\end{table}

Le faible écart-type confirme une performance stable entre les plis, ce qui soutient
la généralisation au-delà de la partition train/test spécifique.

\subsection{Stabilité par graine aléatoire}

Pour vérifier que les résultats ne dépendent pas uniquement de la graine 42, nous
avons répété le pipeline complet (re-partitionnement des données, ré-ajustement
des hyperparamètres) pour les graines 43 et 44. Le Tableau~\ref{tab:seeds} rapporte
les résultats.

\begin{table}[H]
  \centering
  \caption{Stabilité par graine : Full(opt) + SVM relancé sur trois partitions indépendantes.}
  \label{tab:seeds}
  \begin{tabular}{@{}ccc@{}}
    \toprule
    \textbf{Graine} & \textbf{Précision} & \textbf{F1-macro} \\
    \midrule
    42 & 0.9500 & 0.9498 \\
    43 & 0.9100 & 0.9099 \\
    44 & 0.9133 & 0.9130 \\
    \midrule
    \textbf{Moyenne $\pm$ IC 95\,\%} & --- & \textbf{\SeedFOneMean{} $\pm$ \SeedFOneCI} \\
    \bottomrule
  \end{tabular}
\end{table}

Les performances varient d'environ $\pm 0{,}02$ selon la graine, indiquant une robustesse
au partitionnement des données. Cette variation provient des frontières de classe
différentes entre partitions.

\subsection{Intervalles de confiance bootstrap}

À partir de $N = 2000$ rééchantillonnages bootstrap de l'ensemble de test tenu à l'écart,
nous avons estimé des intervalles de confiance à 95\,\% pour le meilleur pipeline
(Tableau~\ref{tab:ci}).

\begin{table}[H]
  \centering
  \caption{Intervalles de confiance bootstrap à 95\,\% pour Full(opt) + SVM (test tenu à l'écart).}
  \label{tab:ci}
  \begin{tabular}{@{}lcc@{}}
    \toprule
    \textbf{Métrique} & \textbf{Estimation ponctuelle} & \textbf{IC 95\,\%} \\
    \midrule
    F1-macro & \BestFOne & [\BestFOneCILow, \BestFOneCIHigh] \\
    Précision & \BestAcc  & [\BestAccCILow, \BestAccCIHigh] \\
    \bottomrule
  \end{tabular}
\end{table}

La Figure~\ref{fig:bootstrap} montre la distribution bootstrap pour le meilleur
jeu de caractéristiques.

\begin{figure}[H]
  \centering
  \includegraphics[width=0.70\linewidth]{bootstrap_ci_best_fs.png}
  \caption{Distribution bootstrap du F1-macro pour Full(opt) + SVM (2000 rééchantillonnages).
           L'IC à 95\,\% est indiqué par des lignes verticales en pointillés.}
  \label{fig:bootstrap}
\end{figure}

\subsection{Tests de significativité}

\subsubsection{Test de McNemar (comparaison par paires sur Full(opt))}

Le test de McNemar compare les schémas d'erreur de deux classifieurs sur le même
ensemble de test. Le Tableau~\ref{tab:mcnemar} résume les comparaisons significatives
avec le SVM.

\begin{table}[H]
  \centering
  \caption{Test de McNemar : SVM contre d'autres classifieurs sur Full(opt). Significatif à $p < 0{,}05$.}
  \label{tab:mcnemar}
  \begin{tabular}{@{}llcc@{}}
    \toprule
    \textbf{Modèle A} & \textbf{Modèle B} & \textbf{$p$-valeur} & \textbf{Significatif ?} \\
    \midrule
    SVM & Random Forest  & \McNemarRF  & Oui \\
    SVM & XGBoost        & \McNemarXGB & Oui \\
    SVM & KNN            & \McNemarKNN & Non \\
    \bottomrule
  \end{tabular}
\end{table}

Le SVM est significativement meilleur que les ensembles d'arbres, mais pas
significativement meilleur que le KNN --- en cohérence avec le classement où
les deux figurent parmi les meilleures configurations.

\subsubsection{Test de Friedman}

Un test de Friedman sur tous les classifieurs au sein du jeu Full(opt) a donné
une statistique $\chi^2$ correspondant à $p = \FriedmanFullP$, confirmant que
le choix du classifieur a un effet statistiquement significatif sur les performances
pour ce jeu de caractéristiques.

La Figure~\ref{fig:statistical_summary} fournit une synthèse visuelle des résultats
de validation.

\begin{figure}[H]
  \centering
  \includegraphics[width=0.80\linewidth]{statistical_validation_summary.png}
  \caption{Synthèse de la validation statistique : IC bootstrap, résultats des plis
           de validation croisée et tests de significativité pour la meilleure configuration.}
  \label{fig:statistical_summary}
\end{figure}
